\documentclass[10pt,a4paper]{amsart}
\usepackage[margin=19mm]{geometry}
\usepackage{amsmath,amssymb,mathtools}
\usepackage[scr=rsfs,cal=euler]{mathalfa}
\usepackage{xfrac}
\usepackage[unicode,hidelinks,bookmarks=false]{hyperref}
\newcommand{\ie}{i.e.\ }
\newcommand{\cf}{cf.\ }
\newcommand{\resp}{respectively\ }
\newcommand{\Subsection}{\subsection}
\providecommand{\nobreakdash}{\nobreak\hbox{-}}
\setlength{\emergencystretch}{2em}
\multlinegap0pt
\usepackage{amssymb}
\newtheorem{thm}{Theorem}
\newtheorem{prop}{Proposition}
\def\Hess{{\rm Hess}}
\def\Rc{\mbox{Rc}}
\def\Scal{{\rm Scal}}
\def\W{{\mathcal W}}
\def\p{\partial}
\title{A note on the equality case of Bray's conjecture}
\author{Xu-Qian Fan, Chengjie Yu$^\dagger$}
\date{Source: arXiv:2608.22182v1 (CC BY 4.0)}
\begin{document}
\maketitle

\noindent\emph{Source and licence.} A note on the equality case of Bray's conjecture. Version arXiv:2608.22182v1 (CC BY 4.0), distributed by arXiv under the Creative Commons Attribution 4.0 International licence, http://creativecommons.org/licenses/by/4.0/, as recorded in the arXiv licence metadata for this exact version and shown on the abstract page https://arxiv.org/abs/2608.22182v1. Full licence notice: \texttt{licenses/arxiv-2608.22182-CC-BY-4.0.txt}.\par\smallskip

\noindent\textit{Editorial scope.} Counted unit: the article's prop prop-derivative-nu (source0.tex lines 224--234). The article's complete original proof of it is delivered with it (source0.tex lines 235--246, one \texttt{proof} environment of 12 lines). No mathematics is written, repaired or re-derived: the statement and the proof are the article's own lines, in the article's own notation, and no retained line is substituted or re-flowed.  Retained uncounted as the source-local context the proof needs: lines 203--214.  Nothing was omitted from any retained span, so the package carries no omission record.  No retained line contains a citation, so the wrapper carries no bibliography and no bibliography entry.  No retained line contains a figure environment, an image-inclusion command or drawing code, so no image is delivered and the package declares no asset dependency.  Editorial formatting only, in addition: an AMS \texttt{amsart} preamble that restates the article's own declarations and macros used by the retained lines, namely the environments \texttt{thm}, \texttt{prop} on separate counters, together with the standard packages those lines need.  The retained environments print in this document as Theorem 1, Proposition 1 in order of appearance, whereas the article prints the same items with the numbering of its own full document.  The complete native source of the article as distributed in the arXiv e-print for this exact version is bundled unchanged as \texttt{sources/208311/source0.tex}.
\begin{thm}\label{thm-derivative-W}
Let $M^n$ be a closed manifold and $g(t)$ be a Ricci flow on $M$. Let $\tau(t)$ be a positive function and $f(t)$ be a family of smooth functions on $M$ such that $\frac{d\tau}{dt}=-1$ and $u(x,t):=(4\pi\tau(t))^{-\frac{n}{2}}e^{-f(x,t)}$ satisfies the conjugate heat equation:
$$\frac{\p u}{\p t}=-\Delta_{g(t)} u+\Scal_{g(t)}u.$$
Then,
\begin{equation}\label{eq-derivative-W}
\begin{split}
&\frac{d}{dt}\mathcal W(g(t),f(t),\tau(t))\\
=&2\tau(t)\int_M\left\|\Rc_{g(t)}+\Hess_{g(t)}(f(t))-\frac{1}{2\tau(t)}g(t)\right\|_{g(t)}^2(4\pi\tau(t))^{-\frac{n}{2}}e^{-f(t)}dV_{g(t)}
\end{split}
\end{equation}
where $\Delta_g$ and $\Hess_g$ are the Laplacian operator and the Hessian operator w.r.t. $g$ and $\|\cdot\|_g$ means taking norm w.r.t. $g$.
\end{thm}

\begin{prop}\label{prop-derivative-nu}
Let $M^n$ be a closed manifold and $g(t)$ with $t\in [a,b]$ be a Ricci flow on $M$. Then, for any $t_0\in (a,b]$, 
\begin{equation}\label{eq-derivative-nu}
D_-\nu(t_0)\geq 2\tau_0\int_M\left\|\Rc_g(t_0)+\Hess_{g(t_0)}(f_0)-\frac{1}{2\tau_0}g(t_0)\right\|_{g(t)}^2(4\pi \tau_0)^{-\frac{n}{2}}e^{-f_0}dV_{g(t_0)}
\end{equation}
where $(f_0,\tau_0)$ is the minimizer so that 
$$\nu(g(t_0))=\W(g(t_0),f_0,\tau_0).$$
Here $\nu(t):=\nu(g(t))$ and 
$$D_-\nu(t_0)=\liminf_{t\to t_0^-}\frac{\nu(t)-\nu(t_0)}{t-t_0}.$$
In particular, if there is a $t_0\in (a,b]$ such that $\nu'(t_0)=0$, then $g(t)$ is a gradient shrinking soliton. 
\end{prop}

\begin{proof}
Let $u_0=(4\pi \tau_0)^{-\frac{n}{2}}e^{-f_0}$, $u(x,t)$ be the solution to the conjugate heat equation:
$$\frac{\p u}{\p t}=-\Delta_{g(t)} u+\Scal_{g(t)} u$$
with $u(t_0)=u_0$ and $$\tau(t)=\tau_0+t_0-t.$$
Let $f(x,t)$ be such that $u(x,t)=(4\pi \tau(t))^{-\frac{n}{2}}e^{-f(x,t)}$. Then
$$\W(g(t),f(t),\tau(t))\geq\nu(t)$$
for $t<t_0$, by the definition of $\nu$, and 
$$\W(g(t_0),f(t_0),\tau(t_0))=\nu(t_0)$$
by that $f(t_0)=f_0$ and $\tau(t_0)=\tau_0$. So, 
$$\frac{\W(g(t),f(t),\tau(t))-\W(g(t_0),f(t_0),\tau(t_0))}{t-t_0}\leq \frac{\nu(t)-\nu(t_0)}{t-t_0}$$
for $t<t_0$. Applying $\displaystyle\liminf_{t\to t_0^-}$ to the both sides of the inequality above and using Theorem \ref{thm-derivative-W}, we get \eqref{eq-derivative-nu}. This completes the proof of the proposition.
\end{proof}

\end{document}
